\chapter{Zusammenfassung und Ausblick}

\section{Zusammenfassung}
Die vorliegende Diplomarbeit befasste sich mit der Konzeption und Implementierung von "Tapyre", einem KI-integrierten Produktivitätssystem, das aus zwei zentralen Komponenten besteht: einem Agentic-AI-Framework für die Desktop-Automatisierung und einer semantischen Suchplattform für wissenschaftliche Publikationen. Ziel war es, moderne Ansätze der künstlichen Intelligenz, insbesondere Große Sprachmodelle (LLMs), in praxisnahe Werkzeuge zu überführen, die den alltäglichen Arbeitsablauf von Entwicklern und Forschern unterstützen.

Die Entwicklung des Agentic-AI-Frameworks basierte auf einer Plugin-Architektur, die es dem System ermöglicht, dynamisch um neue Fähigkeiten erweitert zu werden. Es wurde ein Agent entwickelt, der in der Lage ist, Anweisungen in natürlicher Sprache zu interpretieren und diese mithilfe der zur Verfügung stehenden Plugins in konkrete Aktionen umzusetzen. Die technische Umsetzung erfolgte in Python, um eine enge Anbindung an gängige KI-Bibliotheken zu gewährleisten.

Parallel dazu wurde das System Tapyre Paper Search als Backend für die Verarbeitung und semantische Suche von Forschungsarbeiten implementiert. Die Architektur setzt auf einen modularen Aufbau, der die Datenakquise (z.B. von arXiv), die PDF-Verarbeitung, die Erzeugung von Vektor-Embeddings mit dem domänenspezifischen Modell Specter2 sowie die Speicherung in einer relationalen Datenbank (MySQL) für Metadaten und einer Vektordatenbank (Qdrant) für die semantische Suche klar voneinander trennt. Diese Pipeline bildet das Fundament für die effiziente und skalierbare Erschließung großer Mengen wissenschaftlicher Literatur.

Ein wesentlicher Fokus der Arbeit lag auf der Implementierung von zwei unterschiedlichen Benutzeroberflächen, die auf die jeweiligen Anwendungsfälle zugeschnitten sind:
\begin{itemize}
	\item Die \textbf{Tapyre Desktop-Anwendung}, eine native GTK-Anwendung in Python, die als schlankes Overlay fungiert und eine schnelle, kommandozeilenähnliche Interaktion mit dem KI-Agenten ermöglicht. Der Architekturansatz legte Wert auf eine tiefe Systemintegration und eine performante, non-blocking Ausführung von Agenten-Aufgaben durch den Einsatz von Threading.
	\item Die \textbf{Tapyre Paper Search Webanwendung}, eine moderne Web-App auf Basis von Next.js, React und TypeScript. Sie bietet eine zugängliche, browserbasierte Schnittstelle für die semantische Suche. Die Architektur folgt etablierten Mustern wie der komponenten-basierten Entwicklung und dem unidirektionalen Datenfluss, um eine wartbare und erweiterbare Codebasis zu schaffen.
\end{itemize}
Im Zuge der Implementierung wurden umfassende architektonische Entscheidungen getroffen, alternative Technologien evaluiert und die gewählten Stacks detailliert umgesetzt, von der Backend-Logik bis hin zum responsiven Frontend-Design.

\section{Reflexion und Ergebnisse}
Das Projekt hat erfolgreich demonstriert, dass es möglich ist, mit modernen KI-Technologien und durchdachten Softwarearchitekturen leistungsfähige und zugleich ressourceneffiziente Werkzeuge zu entwickeln.
\begin{itemize}
	\item Es wurde ein funktionsfähiges Agentic-AI-Framework mit einem erweiterbaren Plugin-System geschaffen.
	\item Die Paper-Search-Pipeline ist in der Lage, wissenschaftliche Publikationen zu verarbeiten, zu indexieren und sowohl semantisch als auch über Metadaten durchsuchbar zu machen.
	\item Die beiden entwickelten Benutzeroberflächen bieten jeweils eine intuitive und auf den Anwendungsfall optimierte Interaktionsmöglichkeit mit dem System.
\end{itemize}
Die Evaluation alternativer Technologien sowie die detaillierte Ausarbeitung der Systemarchitekturen haben zudem wertvolle Erkenntnisse über die jeweiligen Vor- und Nachteile der verschiedenen Ansätze (z.B. native Desktop-Entwicklung vs. Web-Frontend) geliefert. Die Analyse des Betriebs hat gezeigt, dass die Systeme grundsätzlich stabil laufen, aber auch Limitationen aufweisen, insbesondere im Bereich der Skalierbarkeit der semantischen Suche unter begrenzten Hardwareressourcen.

\section{Ausblick}
Das Tapyre-System bietet eine solide Grundlage für zahlreiche zukünftige Erweiterungen und Forschungsrichtungen. Die modulare Architektur beider Teilsysteme erleichtert die Weiterentwicklung in verschiedenen Bereichen.

\paragraph{Erweiterung des Tapyre-Agenten}
Das Plugin-System des Desktop-Agenten ist die offensichtlichste Anlaufstelle für Erweiterungen. Zukünftige Arbeiten könnten sich auf die Entwicklung und Integration weiterer Plugins konzentrieren, wie zum Beispiel:
\begin{itemize}
	\item Ein \textbf{Kalender-Plugin} zur Verwaltung von Terminen.
	\item Ein \textbf{E-Mail-Plugin} zum Lesen und Verfassen von E-Mails.
	\item Ein \textbf{Dateisystem-Plugin} für komplexe Dateioperationen.
	\item Eine tiefere Integration mit der Tapyre Paper Search API, sodass der Agent selbstständig Literaturrecherchen durchführen kann.
\end{itemize}

\paragraph{Verbesserungen für Tapyre Paper Search}
Auch die Paper-Search-Plattform bietet zahlreiche Möglichkeiten für zukünftige Verbesserungen:
\begin{itemize}
	\item \textbf{Integration weiterer Datenquellen:} Neben arXiv könnten weitere Plattformen wie PubMed, ACM Digital Library oder Google Scholar angebunden werden, um den durchsuchbaren Datenbestand zu vergrößern.
	\item \textbf{Verbesserte UI/UX:} Die Benutzeroberfläche könnte um weitere Features wie eine Suchhistorie, das Speichern von Suchen, personalisierte Empfehlungen oder die Visualisierung von Autorennetzwerken erweitert werden.
	\item \textbf{Hybride Suchansätze:} Die Kombination aus semantischer Suche und BM25-Keyword-Suche könnte weiter verfeinert werden, um die Suchrelevanz zu maximieren.
\end{itemize}

\paragraph{Perspektive eines KI-gestützten Forschungsassistenten}
Im Zuge der Entwicklung dieser Diplomarbeit entstand zudem eine weiterführende Fragestellung: Wie viel strukturiertes Wissen lässt sich tatsächlich aus großen Sammlungen wissenschaftlicher Publikationen extrahieren, und in welchem Ausmaß könnten KI-Systeme zukünftige Forschungsprozesse unterstützen?

Die Kombination aus einem umfangreichen Paper-RAG-System und agentischen KI-Ansätzen eröffnet hierbei interessante Perspektiven. Ein autonomer Research-Agent könnte beispielsweise relevante Publikationen zu einem bestimmten Thema identifizieren, deren Inhalte analysieren und daraus strukturierte Zusammenfassungen oder Literaturüberblicke erstellen. Dadurch könnten insbesondere zeitaufwendige Rechercheprozesse teilweise automatisiert werden.

Darüber hinaus wäre denkbar, dass ein solcher Agent aus der analysierten Literatur selbstständig neue Forschungsfragen ableitet und diese schrittweise bearbeitet. Beispielsweise könnte der Agent offene Fragestellungen innerhalb eines Themenbereichs identifizieren, zusätzliche Literatur recherchieren, verschiedene Ansätze vergleichen und daraus Hypothesen oder mögliche Forschungsrichtungen formulieren.

Langfristig könnte ein solches System nicht nur einzelne Papers analysieren, sondern auch Zusammenhänge zwischen verschiedenen Arbeiten erkennen, Forschungstrends identifizieren und neue Ideen oder Forschungsrichtungen aufzeigen. Obwohl solche Systeme derzeit noch mit erheblichen technischen und methodischen Herausforderungen verbunden sind, stellt die Kombination aus semantischer Suche, großen Sprachmodellen und agentischen Systemen eine vielversprechende Grundlage für zukünftige Entwicklungen dar.